\chapter*{Exposé XXI. Root Data}
\addcontentsline{toc}{chapter}{Exposé XXI. Root Data}

\noindent
This Exposé\footnote{Editors' note: Version of 13 October 2024.}
collects, in the absence of a suitable reference,\footnote{%
Editors' note: For results on root systems (Sections 1--5), one may
consult \emph{Groupes et algèbres de Lie}, Chapter VI.}
known results on root data (that is, ``abstract'' root systems), most of
which will be used below.

\medskip
\noindent\textbf{Notation.}
We denote by \(\Qp\) the set of positive (or zero) rational numbers;
\(\mathbb Z\cap\Qp=\mathbb N\). Let \(V\) be a \(\mathbb Q\)-vector
space. If \(A\) (respectively \(B\)) is a subset of \(\mathbb Q\)
(respectively of \(V\)), we denote by \(A\cdot B\) the image of
\(A\otimes B\) under the morphism
\(\mathbb Q\otimes V\to V\), in other words the set of linear
combinations of elements of \(B\) with coefficients in \(A\). We write
\(-B=\{-1\}B\). We denote by \(E-F\) the set of elements of \(E\)
that do not belong to \(F\).

\section{Generalities}

\SGAInlineRefTarget{sga3:xxi:subsection:1-1}
\subsection{Definitions and first properties}

\begin{definition}\label{sga3:XXI:1.1.1}
Let \(M\) and \(M^\ast\) be two free finite-type \(\mathbb Z\)-modules
in duality. Write
\[
  V=M\otimes\mathbb Q,\qquad V^\ast=M^\ast\otimes\mathbb Q;
\]
these are two \(\mathbb Q\)-vector spaces in duality. We identify \(M\)
(respectively \(M^\ast\)) with a subset of \(V\) (respectively
\(V^\ast\)). The canonical bilinear form on
\(M^\ast\times M\) (respectively \(V^\ast\times V\)) is denoted
\((\, ,\,)\).

Let \(R\) be a finite subset of \(M\). Suppose given a map
\(\alpha\mapsto\alpha^\ast\) from \(R\) to \(M^\ast\); the set of
\(\alpha^\ast\), for \(\alpha\in R\), is denoted \(R^\ast\).
To every \(\alpha\in R\) associate the endomorphism \(s_\alpha\)
(respectively \(s_\alpha^\ast\)) of \(M\) and \(V\) (respectively of
\(M^\ast\) and \(V^\ast\)) given by
\begin{align}
 s_\alpha(x)&=x-(\alpha^\ast,x)\alpha,
 &s_\alpha&=\operatorname{id}-\alpha^\ast\otimes\alpha,
 \tag{1}\label{sga3:XXI:eq:1}\\
 s_\alpha^\ast(u)&=u-(u,\alpha)\alpha^\ast,
 &s_\alpha^\ast&=\operatorname{id}-\alpha\otimes\alpha^\ast.
 \tag{1${}^\ast$}\label{sga3:XXI:eq:1star}
\end{align}
The pair \((R,R^\ast)\) (more precisely, the pair
\((R,R\to M^\ast)\)) is called a \emph{root datum} in
\((M,M^\ast)\), and \((M,M^\ast,R,R^\ast)\) is called a root datum,
if the following axioms hold:
\begin{align*}
 (\mathrm{DR\,I})\quad&
    (\alpha^\ast,\alpha)=2
    &&\text{for every }\alpha\in R,\\
 (\mathrm{DR\,II})\quad&
    s_\alpha(R)\subset R,\qquad
    s_\alpha^\ast(R^\ast)\subset R^\ast
    &&\text{for every }\alpha\in R.
\end{align*}
The set \(R\) is called the \emph{root system} of the root datum
\(\cR=(M,M^\ast,R,R^\ast)\). The elements of \(R\) (respectively
\(R^\ast\)) are called the \emph{roots} (respectively
\emph{coroots}) of the root datum.
\end{definition}

\begin{remark}\label{sga3:XXI:1.1.2}
Axiom \((\mathrm{DR\,I})\) is equivalent to any one of the following
properties:
\begin{align}
 s_\alpha s_\alpha&=\operatorname{id},
 &s_\alpha^\ast s_\alpha^\ast&=\operatorname{id},
 \tag{2,2${}^\ast$}\label{sga3:XXI:eq:2}\\
 s_\alpha(\alpha)&=-\alpha,
 &s_\alpha^\ast(\alpha^\ast)&=-\alpha^\ast.
 \tag{3,3${}^\ast$}\label{sga3:XXI:eq:3}
\end{align}
\end{remark}

\begin{remark}\label{sga3:XXI:1.1.3}
Axioms \((\mathrm{DR\,I})\) and \((\mathrm{DR\,II})\) imply
\[
 R=-R,\qquad R^\ast=-R^\ast,\qquad
 0\notin R,\qquad 0\notin R^\ast.
\]
\end{remark}

\begin{lemma}\label{sga3:XXI:1.1.4}
The map \(R\to R^\ast\) is a bijection. More generally, if
\(\alpha,\beta\in R\) and
\((\alpha^\ast,x)=(\beta^\ast,x)\) for every \(x\in R\), then
\(\alpha=\beta\).
\end{lemma}

\begin{proof}
Indeed, one then has
\[
 s_\beta(\alpha)=\alpha-2\beta,\qquad
 s_\alpha(\beta)=\beta-2\alpha.
\]
It follows at once that
\[
 s_\beta s_\alpha(\alpha)
   =2\beta-\alpha=\alpha+2(\beta-\alpha),\qquad
 s_\beta s_\alpha(\beta-\alpha)
   =s_\beta(\beta-\alpha)=\beta-\alpha.
\]
Hence, by \((\mathrm{DR\,II})\),
\[
 (s_\beta s_\alpha)^n(\alpha)
   =\alpha+2n(\beta-\alpha)\in R.
\]
Since \(R\) is finite, \(\beta-\alpha=0\).
\end{proof}

\begin{corollary}\label{sga3:XXI:1.1.5}
The inverse map \(R^\ast\to R\) defines a root datum
\[
 \cR^\ast=(M^\ast,M,R^\ast,R),
\]
called the \emph{dual} of \(\cR\).\footnote{Editors' note:
\((\alpha^\ast)^\ast=\alpha\).}
\end{corollary}

\begin{definition}\label{sga3:XXI:1.1.6}
We denote by \(\Gamma_0(R)\) the subgroup of \(M\) generated by \(R\).
We denote by \(\VR(R)\) the vector subspace of \(V\) generated by \(R\),
that is,
\[
 \VR(R)=\Gamma_0(R)\otimes\mathbb Q.
\]
Applying these definitions to \(\cR^\ast\), one likewise obtains
\(\Gamma_0(R^\ast)\) and \(\VR(R^\ast)\).

The \emph{reductive rank} of \(\cR\) is the number
\[
 \rgred(\cR)
 =\rank(M)=\dim(V)=\dim(V^\ast)=\rank(M^\ast)=\rgred(\cR^\ast).
\]
The \emph{semisimple rank} of \(\cR\) is the number
\[
 \rgss(\cR)
 =\rank(R)=\rank(\Gamma_0(R))=\dim(\VR(R)).
\]
Thus \(\rgss(\cR)\leq\rgred(\cR)\).

We shall see below that
\(\rgss(\cR)=\rgss(\cR^\ast)\), that is, that \(\VR(R)\) and
\(\VR(R^\ast)\) have the same dimension.
\end{definition}

\begin{definition}\label{sga3:XXI:1.1.7}\leavevmode\par
The root datum \(\cR\) is called \emph{semisimple} (respectively
\emph{trivial}) if
\(\rgss(\cR)=\rgred(\cR)\) (respectively \(\rgss(\cR)=0\)).
Thus \(\cR\) is trivial if and only if \(R\) is empty. The trivial root
datum of reductive rank zero is denoted
\[
 0=(\{0\},\{0\},\varnothing,\varnothing).
\]
\end{definition}

\begin{definition}\label{sga3:XXI:1.1.8}
We denote by \(W(\cR)\) the group of transformations of \(M\) generated
by the \(s_\alpha\), \(\alpha\in R\). It is called the \emph{Weyl
group} of \(\cR\). We write
\[
 W^\ast(\cR)=W(\cR^\ast).
\]
Then \(W(\cR)\) acts on \(R\), \(\Gamma_0(R)\), \(\VR(R)\), \(M\),
and \(V\). If \(w\in W(\cR)\) and \(x\in M\) (respectively
\(x\in V\)), then
\[
 wx-x\in\Gamma_0(R)
 \quad\text{(respectively }wx-x\in\VR(R)\text{)};
\]
this follows immediately from formula \eqref{sga3:XXI:eq:1}.
Likewise for \(W^\ast(\cR)\).
\end{definition}

\begin{lemma}\label{sga3:XXI:1.1.9}
For every \(\alpha\in R\), \(x\in V\), and \(u\in V^\ast\), one has
\begin{equation}
 (s_\alpha^\ast(u),s_\alpha(x))=(u,x).
 \tag{4}\label{sga3:XXI:eq:4}
\end{equation}
\end{lemma}

\begin{proof}
Taking the scalar product of \eqref{sga3:XXI:eq:1} and
\eqref{sga3:XXI:eq:1star}, the left-hand side is
\[
 (u,x)+(u,\alpha)(\alpha^\ast,x)
       \bigl((\alpha^\ast,\alpha)-2\bigr)=(u,x).
\]
\end{proof}

\begin{remark}\label{sga3:XXI:1.1.10}
If one assumes \(0\notin R\) and \(0\notin R^\ast\), then
Lemma~\ref{sga3:XXI:1.1.9} is equivalent to
\((\mathrm{DR\,I})\).
\end{remark}

\begin{corollary}\label{sga3:XXI:1.1.11}
Let \(h\mapsto h^\vee\) denote the isomorphism from
\(\GL(M)\) onto \(\GL(M^\ast)\) that associates to \(h\) its
contragredient. Formula \eqref{sga3:XXI:eq:4} may then also be written
\begin{equation}
 (s_\alpha)^\vee=s_\alpha^\ast.
 \tag{5}\label{sga3:XXI:eq:5}
\end{equation}
\end{corollary}

\begin{corollary}\label{sga3:XXI:1.1.12}
The preceding isomorphism induces an isomorphism from \(W(\cR)\) onto
\(W^\ast(\cR)\).
\end{corollary}

\begin{scholium}\label{sga3:XXI:1.1.13}
By the preceding result, we shall identify \(W\) and \(W^\ast\), and
regard \(W\) as a group of transformations of
\[
 R,\ R^\ast,\ M,\ M^\ast,\ \Gamma_0(R),\ \Gamma_0(R^\ast),\
 V,\ V^\ast,\ \VR(R),\ \VR(R^\ast).
\]
We shall write \(s_\alpha\) for \(s_\alpha^\ast\).
\end{scholium}
